\chapter{Methodology}\label{ch:method}

\section{Method}\label{sec:method}

\subsection{Datasets, models, and protocol}\label{sec:method-data}
Three public, de-identified tabular datasets: PIMA Indian Diabetes
(768 rows, 8 features, 34.9\,\% positive)~\cite{smith1988using}, the Cleveland subset of
UCI Heart Disease (303, 13, 45.9\,\%)~\cite{detrano1989international}, and UCI Chronic
Kidney Disease (400, 24, 62.5\,\%)~\cite{rubini2015ckd}; all are redistributable under CC0
or CC BY 4.0 and are vendored with verified checksums. They span a \emph{difficulty gradient} ---
nested-CV ROC-AUC of \AucPima{}, \AucHeart{}, and \AucCkd{} --- because a single-dataset
study cannot ask whether agreement depends on how hard the boundary is. The diseases are
benchmarks; no medical claim is made. In PIMA, biologically impossible zeros in glucose,
blood pressure, skin thickness, insulin, and BMI are recoded as missing while legitimate
zeros in \emph{pregnancies} are retained; the Heart target is binarised at $>0$ with its
nominal codes carried as categories, so the linear model cannot read an ordering into
them.


Six classifiers span linear, single-tree, bagged, boosted, and stacked families: logistic
regression, a decision tree, a random forest, XGBoost~\cite{chen2016xgboost},
LightGBM~\cite{ke2017lightgbm}, and a stacking ensemble~\cite{wolpert1992stacked}.
Stacking is included because it admits no tree-specific explainer, making it the one family
where \shap{}'s own re-run variance is a non-trivial measurement.

Performance and explanation analysis run on separated tracks. \textbf{Track A} estimates
performance by nested cross-validation, five outer folds for estimation and three inner for
tuning, with every transformer inside the pipeline so imputation and scaling are fitted on
inner-training rows only. \textbf{Track B} takes those hyperparameters, refits once on a
stratified 80/20 split, and runs every explainer against that one model per cell. The
separation prevents a specific leak: the partition the explainers see must never have
participated in tuning. We apply \emph{no} synthetic resampling --- imbalance is mild, and
oversampling would alter the data manifold that \lime{}'s sampler and \shap{}'s background
both draw from, confounding the comparison under study.

\subsection{Explainers and metrics}\label{sec:method-metrics}
\shap{} uses \textsc{TreeExplainer} for the four tree models and a model-agnostic
permutation sampler for logistic regression and stacking; \lime{} uses
\code{LimeTabularExplainer} with one-hot columns declared categorical. A uniformly random ranking is the control.

\subsection{The hybrid consensus method}\label{sec:method-hybrid}
Both explainers produce a ranking over the same $d$ original features, so they can be
combined at the level of ranks without either having to be rescaled. We evaluate three
fusion rules (Algorithm~\ref{alg:fusion}).

\begin{algorithm}[t]
\caption{Rank-level fusion of two explanations}\label{alg:fusion}
\KwIn{attributions $\phi^{A}, \phi^{B} \in \mathbb{R}^{d}$; rule; weight $w$}
\KwOut{fused ranking over the $d$ original features}
$r^{A} \gets$ normalised ranks of $|\phi^{A}|$\;
$r^{B} \gets$ normalised ranks of $|\phi^{B}|$\;
\Switch{rule}{
  \Case{mean-rank (baseline)}{$r \gets \tfrac{1}{2}(r^{A} + r^{B})$\;}
  \Case{RRF (baseline)}{$s_j \gets \tfrac{1}{60+r^{A}_j} + \tfrac{1}{60+r^{B}_j}$;
        $r \gets$ ranks of $-s$\;}
  \Case{fidelity-weighted (ours)}{$r \gets w\,r^{A} + (1-w)\,r^{B}$\;}
}
\Return{$\operatorname{argsort}(r)$}\;
\end{algorithm}

Two are published baselines: \emph{mean-rank}
fusion, which averages the normalised ranks, and \emph{reciprocal-rank fusion}, standard in
information retrieval, where each list contributes $1/(60+r_j)$ so one badly wrong ranking
degrades the result less than an arithmetic mean would~\cite{cormack2009reciprocal}.

The third rule is ours. \emph{Fidelity-weighted fusion} does not treat the two explainers
as equally trustworthy. It weights them by the faithfulness each actually demonstrates:
$w = \aopc{}^{\text{cal}}_{\shap} / (\aopc{}^{\text{cal}}_{\shap} +
\aopc{}^{\text{cal}}_{\lime})$, so the explainer whose rankings move the model's
prediction more on held-out instances carries more of the fused rank. The weight is
estimated \emph{once per cell} on a calibration split of 30\,\% of the explained
instances, and \RQ{6} is then tested on the remaining 70\,\% only. That disjointness is
the point: a weight fitted on the same instances the hypothesis is tested on would make the
rule win by construction, and the split is recorded per cell so the separation is auditable
rather than asserted.

The bar for all three rules is deliberately strict. A fused ranking must beat the
\emph{better} of \shap{} and \lime{} on the same cell --- not their average, which any
sensible combination would beat, and not the weaker of the two.

Let $\phi(x) \in \mathbb{R}^{d}$ be an attribution vector over original features and
$\pi(x)$ the ranking by $|\phi(x)|$. \emph{Agreement} (\RQ{1}) is Spearman $\rho$, Kendall
$\tau_b$, top-$k$ Jaccard, and rank-biased overlap between $\pi_{\shap}$ and
$\pi_{\lime}$. Because both rankings permute the \emph{same} feature set, \rbo{}'s
extrapolation term is a hard floor scaling with $d$ (\RboFloorEight{} at $d=8$ but
\RboFloorTwentyFour{} at $d=24$), so we also report a floor-normalised variant comparable
across datasets. \emph{Stability} (\RQ{2}) is decomposed: re-run variance repeats an
explainer $R=10$ times under different seeds; local robustness draws $m=5$ Gaussian
neighbours at $0.1\sigma$ and estimates $\hat{L}(x)=\max_{x'} \lVert
\phi(x)-\phi(x')\rVert_2 / \lVert x-x'\rVert_2$~\cite{alvarezmelis2018robustness};
retraining stability refits each ensemble on $B=10$ bootstrap resamples and reports
cross-refit rank agreement. \emph{Fidelity} (\RQ{3}) deletes features in ranked order,
replacing them with the training median, and records the fall in the model's probability
for \emph{its own} predicted class: $\aopc{} =
\frac{1}{K}\sum_{k=1}^{K}[f_c(x)-f_c(x_{\setminus \text{top-}k})]$, $K=\min(d,10)$, with
insertion AUC, comprehensiveness, and sufficiency
alongside~\cite{deyoung2020eraser} and a marginal-draw baseline as a sensitivity check.
\emph{Cost} (\RQ{4}) is single-threaded wall-clock seconds around the explainer call.
\emph{Plausibility} (\RQ{5}) is top-5 overlap with a literature-derived reference set
frozen before any explanation was computed.

\subsection{Four traps this comparison had to avoid}\label{sec:pitfalls}
Each of the following silently invalidates a \shap{}-versus-\lime{} comparison, and none
is visible in the output when it goes wrong.

\textbf{The two methods must explain the same quantity.} \lime{} explains predicted
probability; \textsc{TreeExplainer} on boosted models returns log-odds margins by default,
so a direct comparison contrasts different quantities on different scales. We require
probability-space attributions and \emph{verify} it, by checking that
$\text{base}+\sum_j \phi_j$ reproduces \code{predict\_proba}. Median additivity error
across cells is \AdditivityMedianMax{}. In \AdditivityGateFailCells{} of \NumCells{} cells
the pre-registered gate did not pass (PIMA/LightGBM: \AdditivityWorstOutliers{} of 150
instances above the $5\times10^{-3}$ tolerance, worst \AdditivityWorstMax{}, against a
cell median of $1.1\times10^{-9}$); Section~\ref{sec:limitations} states how that is
handled.

\textbf{The reference distribution must be real data.} Summarising the \shap{} background
by $k$-means yields synthetic points, because averaging one-hot columns produces
fractional values the model has never seen. Interventional \textsc{TreeSHAP} evaluated
against such off-manifold references violates its own additivity --- we measured
$9.9\times10^{-3}$ on Heart/LightGBM against $7\times10^{-9}$ with real training rows. The
background is therefore a seeded, class-stratified subsample of actual rows, which also
keeps tree \shap{} deterministic.

\textbf{One-hot encoding fragments a feature across several model inputs.} Left
unaggregated, rankings from different models would be over different vocabularies, and
deleting ``the top five features'' would delete five \emph{columns} --- under-perturbing a
13-feature dataset spread over 25 columns while leaving an 8-feature one untouched.
Attributions are summed over each feature's columns before ranking, and perturbation
operates on whole feature groups.

\textbf{A fidelity score has no scale without a control.} Deleting any few of eight
features moves a prediction, so an \aopc{} of $0.31$ is uninterpretable alone. Every
fidelity number is reported against a random-ranking control.

\subsection{Statistics}\label{sec:method-stats}
Eleven hypotheses were registered in advance with their test, alternative, unit of
analysis, and family (Table~\ref{tab:hypotheses}). All tests are non-parametric;
\emph{p}-values are Holm-corrected within each family over the tests actually performed,
with degenerate cells retained and reported but excluded from the divisor. Every test
carries an effect size and every estimate a 95\,\% bootstrap interval: with hundreds of
instances per cell significance is cheap and the effect size is the finding. \RQ{1} tests
whether concordance exceeds $\rho=0.50$ rather than zero, because testing against zero
would be a straw man; the threshold was fixed before any $\rho$ was computed.


\section{Notation}
\begin{tabular}{@{}ll@{}}
\toprule
Symbol & Meaning \\
\midrule
$d$ & number of original features in a dataset \\
$x$ & one instance, in transformed model-input space \\
$\phi(x) \in \mathbb{R}^{d}$ & attribution vector, aggregated to original features \\
$\pi(x)$ & ranking of features by $|\phi(x)|$, ties broken by index \\
$f_c(\cdot)$ & model probability for the instance's own predicted class $c$ \\
$K$ & deletion depth, $\min(d, 10)$ \\
$R$ & explainer re-runs per instance, 10 \\
$m$ & perturbation neighbours per instance, 5 \\
$B$ & bootstrap model refits, 10 \\
$w$ & weight given to SHAP by the fidelity-weighted fusion rule \\
$\hat{L}(x)$ & local Lipschitz estimate of explanation robustness \\
\bottomrule
\end{tabular}

\section{Hyperparameter search spaces}
\begin{table}[t]
\centering
\caption{Hyperparameter search spaces and SHAP backend assignment. Grids are capped at 12 combinations to fit the compute budget.}
\label{tab:grids}
\small
\begin{tabular}{llrl}
\toprule
Model & Search space & Combinations & SHAP backend \\
\midrule
Logistic Regression & C $\in$ \{0.01, 0.1, 1.0, 10.0\} & 4 & permutation \\
Decision Tree & max\_depth $\in$ \{3, 5, 8\}; min\_samples\_leaf $\in$ \{2, 5, 10\} & 9 & tree \\
Random Forest & max\_depth $\in$ \{unlimited, 8\}; min\_samples\_leaf $\in$ \{1, 2, 5\} & 6 & tree \\
XGBoost & max\_depth $\in$ \{2, 3, 5\}; learning\_rate $\in$ \{0.05, 0.1\} & 6 & tree-prob \\
LightGBM & num\_leaves $\in$ \{7, 15, 31\}; learning\_rate $\in$ \{0.05, 0.1\} & 6 & tree-prob \\
Stacking & final\_estimator\_\_C $\in$ \{0.1, 1.0, 10.0\} & 3 & permutation \\
\bottomrule
\end{tabular}
\end{table}


\section{Complete preprocessing specification}
Numeric columns are median-imputed then standardised; categorical columns are
mode-imputed then one-hot encoded with binary features collapsed to a single column and
unseen categories at transform time encoded as all-zero. Every step is fitted inside a
\texttt{Pipeline}, so under Track A the statistics come from the inner-training rows of
each fold and under Track B from the 80\,\% training partition alone. The mapping from
output columns back to original features is derived from the fitted encoder's own
category lists rather than by parsing generated column names --- the Chronic Kidney
Disease data contains both \texttt{pc} and \texttt{pcc}, and any name-prefix heuristic
assigns one's columns to the other.

\section{Compute budget and profiles}
The pipeline runs under two configuration profiles sharing one code path: a QUICK
profile that validates the entire chain including document builds in about five
minutes, and the FULL profile that produces the reported results. Every figure carries
a visible watermark under any non-reporting profile, so a draft cannot be mistaken for
a result. Recorded settings for the run that produced this document are in
Appendix~\ref{app:environment}.
